\chapter{Fazit und Ausblick}
\label{chap:fazit}

\section{Zusammenfassung der Ergebnisse}
\label{sec:zusammenfassung}

Diese Bachelorarbeit beschäftigt sich mit dem Entwurf und der
prototypischen Entwicklung des Frontends der Webanwendung
\textit{Bazario}. Ziel war eine rollenbasierte Single-Page
Application, über die sich Produkte kaufen und Dienstleistungen buchen
lassen, ohne dass der Nutzer zwischen zwei Anwendungen wechseln muss.

Am Anfang wurden die Motivation und die Anforderungen an das Frontend
behandelt und der Technologiestack darauf aufbauend begründet. Danach
folgten das Konzept mit Komponentenarchitektur, rollenbasiertem
Seitenschutz und \acs{UI}-Gestaltung sowie die Umsetzung ausgewählter
Funktionsbereiche mit React und TypeScript. Die entwickelte Oberfläche
deckt die Funktionsbereiche der drei im Frontend abgebildeten Rollen
ab, vom Katalog über den gemeinsamen Warenkorb für Produkte und
Buchungen bis zum Kontobereich, zum Nachrichtenbereich und zur
Verwaltung von Werbung; sie liegt in Deutsch, Englisch und Arabisch
vor.

Die Bewertung in Kapitel~\ref{chap:testen} erfolgte auf drei Wegen.
Die automatisierten Tests umfassen 36 Testdateien mit 134 Testfällen
sowie vier End-to-End-Tests und laufen fehlerfrei durch. Die
heuristische Evaluation ergab vierzehn Befunde, die sämtlich behoben
wurden. Von den vier geprüften Erfolgskriterien der
\acs{WCAG}~2.1 sind drei erfüllt; nicht erfüllt ist allein der
Kontrast des Fokusrings.

\section{Mögliche Erweiterungen und Ausblick}
\label{sec:ausblick}

Für die Weiterentwicklung des Frontends bieten sich mehrere
Erweiterungen an:

\begin{itemize}
  \item Eine eigene Administrationsoberfläche für Rollenwechselanträge,
  Freigaben, Gebühren und Auszahlungen.
  \item Erweiterte Suchfilter nach Preis, Ort und Verfügbarkeit.
  \item Die Korrektur des Fokusrings sowie eine automatisierte Prüfung
    der Kontrastwerte im Build-Vorgang.
  \item Die Anzeige des Lagerbestands mit einer Warnung im Warenkorb,
  wenn die gewählte Menge den Bestand übersteigt.
  \item Die Ergänzung von Lieferanschrift, Versandkosten und
  Sendungsstatus im Bestellprozess.
  \item Der Ausbau zur progressiven Web-App mit Benachrichtigungen zu
  Bestellungen, Buchungen und Nachrichten.
\end{itemize}
Rein im Frontend umsetzbar sind die Korrektur des Fokusrings sowie die
Administrationsoberfläche, da deren Funktionen bereits über die
Demonstrationsoberfläche bedient werden. Die übrigen Erweiterungen
setzen zusätzliche Schnittstellen im Backend voraus
\cite{alaatki_backend}.

Zusammenfassend stellt das entwickelte Frontend eine tragfähige und
erweiterbare Oberfläche für eine Plattform dar, die den Kauf von
Produkten und die Buchung von Dienstleistungen in einem gemeinsamen
Vorgang zusammenführt.